% Source: linear_time_series_models.tex
\begin{notechapter}
\chapter{Some Linear Time Series Models}\label{ch:linear-time-series}

\section*{Preamble: Key Operators in Time Series Analysis}
Before diving into the models, we introduce two fundamental operators 
used extensively in time series analysis: the \textbf{Backward Shift 
Operator} and the \textbf{Difference Operator}.

\subsection*{Backward Shift Operator}
The \textbf{Backward Shift Operator} $B$ is defined as:
\[
B X_t = X_{t-1}, \quad B^k X_t = X_{t-k}
\]
This operator shifts the time index of a time series backward by one 
or more steps.

\subsection*{Difference Operator}
The \textbf{Difference Operator} $\nabla$ is defined as:
\[
\nabla X_t = X_t - X_{t-1} = (1 - B) X_t
\]
This operator computes the difference between consecutive observations. 
Higher-order differences are defined recursively:
\[
\nabla^d X_t = \nabla (\nabla^{d-1} X_t)
\]
For example, the second difference is:
\[
\nabla^2 X_t = \nabla (X_t - X_{t-1}) = X_t - 2X_{t-1} + X_{t-2} = (1 - B)^2 X_t
\]

\subsection*{Seasonal Difference Operator}
For seasonal time series, the \textbf{Seasonal Difference Operator} 
$\nabla_s$ is defined as:
\[
\nabla_s X_t = X_t - X_{t-s} = (1 - B^s) X_t
\]
where $s$ is the seasonal period (e.g., $s = 12$ for monthly data). 
Higher-order seasonal differences are defined as:
\[
\nabla_s^D X_t = (1 - B^s)^D X_t
\]

\section*{Introduction}
This chapter introduces various probability models for time series, 
focusing on linear models such as autoregressive (AR), moving average 
(MA), and mixed ARMA models. The concept of stationarity is central 
to these models, and we explore how to identify and fit these models 
to time series data.

\section{Stochastic Processes and Their Properties}
A \textbf{stochastic process} is a collection of random variables 
ordered in time. It can be either continuous or discrete. In this 
chapter, we focus on discrete-time processes, where observations 
are made at equal intervals.

\subsection*{Key Properties}
\begin{itemize}
    \item \textbf{Mean function}: $\mu(t) = E[X(t)]$
    \item \textbf{Variance function}: $\sigma^2(t) = \text{Var}[X(t)]$
    \item \textbf{Autocovariance function (acv.f.)}: 
          $\gamma(t_1, t_2) = \text{Cov}[X(t_1), X(t_2)]$
    \item \textbf{Autocorrelation function (ac.f.)}: 
          $\rho(\tau) = \gamma(\tau)/\gamma(0)$
\end{itemize}

\section{Stationary Processes}
A time series is \textbf{strictly stationary} if the joint distribution 
of any set of observations is invariant to time shifts. A process is 
\textbf{second-order stationary} (or weakly stationary) if its mean 
and variance are constant, and the autocovariance depends only on 
the lag $\tau$.

\subsection*{Importance of Stationarity}
Stationarity is crucial because much of the theory of time series 
analysis is based on stationary processes. Non-stationary processes 
can often be transformed into stationary ones by differencing or 
other techniques.

\section{Properties of the Autocorrelation Function (ac.f.)}
The autocorrelation function (ac.f.) is a key tool for analyzing 
time series. It measures the correlation between observations at 
different lags.

\subsection*{Key Properties of the ac.f.}
\begin{itemize}
    \item The ac.f. is an even function: $\rho(\tau) = \rho(-\tau)$.
    \item The ac.f. is bounded: $|\rho(\tau)| \leq 1$.
    \item The ac.f. does not uniquely identify the underlying model, 
          as different processes can have the same ac.f.
\end{itemize}

\section{Purely Random Processes}
A \textbf{purely random process} consists of a sequence of independent 
and identically distributed (i.i.d.) random variables, often assumed 
to be normally distributed with mean zero and variance $\sigma^2_Z$.

\subsection*{Properties}
\begin{itemize}
    \item The ac.f. of a purely random process is zero for all lags 
          except lag 0, where it is 1.
    \item Such processes are also called \textbf{white noise}.
\end{itemize}

\section{Random Walks}
A \textbf{random walk} is a process where each value is the sum of 
the previous value and a purely random error term:
\[
X_t = X_{t-1} + Z_t
\]
The process is customarily started at zero when t = 0, so that
\[
X_1 = Z_1,\; \text{and}\; X_t = \sum_{i=1}^{t}Z_i
\]

Random walks are non-stationary, but their first differences are 
stationary. Using the difference operator, this can be written as:
\[
\nabla X_t = Z_t
\]

\section{Moving Average (MA) Processes}
An \textbf{MA(q) process} is defined as:
\[
X_t = Z_t + \beta_1 Z_{t-1} + \cdots + \beta_q Z_{t-q}
\]
Using the backward shift operator, this can be written as:
\[
X_t = \theta(B) Z_t
\]
where $\theta(B) = 1 + \beta_1 B + \cdots + \beta_q B^q$ is the MA 
polynomial.

\subsection*{Introduction to MA Processes}
The moving average (MA) model is a simple yet powerful tool for 
modeling time series. It predicts future values based on past 
forecast errors, making it particularly useful for capturing 
short-term fluctuations in the data. The order $q$ of the MA model 
determines how many past errors are used for prediction.

\subsection*{Properties}
\begin{itemize}
    \item The ac.f. of an MA(q) process "cuts off" after lag $q$, 
          meaning $\rho(k) = 0$ for $k > q$.
    \item MA processes are always stationary.
    \item The process must satisfy the \textbf{invertibility condition} 
          to ensure uniqueness. This condition requires that the roots 
          of the MA polynomial $\theta(B)$ lie outside the unit circle.
\end{itemize}

\subsection*{Invertibility Condition}
A process is invertible if there exists a polynomial $\theta^{-1}(B)$ 
such that:
\[
\theta^{-1}(B) \theta(B) = 1
\]
This ensures that the MA process can be expressed as an infinite AR 
process, which is useful for forecasting and interpretation.

\section{Autoregressive (AR) Processes}
An \textbf{AR(p) process} is defined as:
\[
X_t = \alpha_1 X_{t-1} + \cdots + \alpha_p X_{t-p} + Z_t
\]
Using the backward shift operator, this can be written as:
\[
\phi(B) X_t = Z_t
\]
where $\phi(B) = 1 - \alpha_1 B - \cdots - \alpha_p B^p$ is the AR 
polynomial.

\subsection*{Introduction to AR Processes}
The autoregressive (AR) model is a fundamental tool in time series 
analysis. It predicts future values based on past values, making it 
particularly useful for capturing trends and dependencies in the data. 
The order $p$ of the AR model determines how many past values are used 
for prediction.

\subsection*{Characteristic Equation}
The \textbf{characteristic equation} of an AR(p) process is given by:
\[
1 - \alpha_1 B - \alpha_2 B^2 - \cdots - \alpha_p B^p = 0
\]
The roots of this equation determine the stationarity of the process.

\subsection*{Stationarity Condition}
An AR(p) process is stationary if all the roots of the characteristic 
equation lie outside the unit circle. This ensures that the process 
does not explode over time.

\subsection*{Autocorrelation Function (ac.f.)}
The ac.f. of an AR(p) process decays exponentially or as a damped 
sinusoid, depending on the roots of the characteristic equation. 
For example:
\begin{itemize}
    \item If the roots are real, the ac.f. decays exponentially.
    \item If the roots are complex, the ac.f. exhibits damped 
          sinusoidal behavior.
\end{itemize}

\subsection*{Yule-Walker Equations}
The parameters of an AR(p) process can be estimated using the 
\textbf{Yule-Walker equations}:
\[
\rho(k) = \alpha_1 \rho(k-1) + \cdots + \alpha_p \rho(k-p), \quad k \geq 1
\]
These equations relate the autocorrelations to the AR coefficients.

\section{Mixed ARMA Models}
An \textbf{ARMA(p, q) model} combines AR and MA components:
\[
X_t = \alpha_1 X_{t-1} + \cdots + \alpha_p X_{t-p} + Z_t + \beta_1 Z_{t-1} 
+ \cdots + \beta_q Z_{t-q}
\]
Using the backward shift operator, this can be written as:
\[
\phi(B) X_t = \theta(B) Z_t
\]
where $\phi(B) = 1 - \alpha_1 B - \cdots - \alpha_p B^p$ is the AR 
polynomial, and $\theta(B) = 1 + \beta_1 B + \cdots + \beta_q B^q$ 
is the MA polynomial.

\subsection*{Introduction to ARMA Processes}
The ARMA model combines the strengths of both AR and MA models, 
allowing for more flexible and accurate predictions. The AR component 
captures the dependency on past values, while the MA component captures 
the dependency on past errors. Together, they provide a powerful tool 
for modeling stationary time series.

\subsection*{Stationarity and Invertibility}
\begin{itemize}
    \item The ARMA process is stationary if the roots of $\phi(B)$ 
          lie outside the unit circle.
    \item The process is invertible if the roots of $\theta(B)$ 
          lie outside the unit circle.
\end{itemize}

\subsection*{Autocorrelation Function (ac.f.)}
The ac.f. of an ARMA process is a mixture of the ac.f.s of the AR 
and MA components. It generally decays exponentially after lag $q$.

\section{Integrated ARMA (ARIMA) Models}
An \textbf{ARIMA(p, d, q) model} is used for non-stationary time series 
that can be made stationary by differencing $d$ times. The model is 
written as:
\[
\phi(B) \nabla^d X_t = \theta(B) Z_t
\]
or equivalently:
\[
\phi(B) (1 - B)^d X_t = \theta(B) Z_t
\]
where $\nabla^d = (1 - B)^d$ is the $d$-th difference operator.

\subsection*{Introduction to ARIMA Processes}
The ARIMA model extends the ARMA model by incorporating differencing 
to handle non-stationary time series. Differencing removes trends and 
makes the series stationary, allowing the ARMA model to be applied. 
The order $d$ of differencing determines how many times the series 
is differenced.

\subsection*{Properties}
\begin{itemize}
    \item The differenced series $\nabla^d X_t$ is an ARMA(p, q) 
          process.
    \item ARIMA models are widely used in practice, especially for 
          economic and financial time series.
    \item The differencing operation removes trends and makes the 
          series stationary.
    \item The parameters $p$, $d$, and $q$ must be chosen carefully 
          to ensure a good fit.
\end{itemize}

\section{Seasonal ARIMA (SARIMA) Models}
A \textbf{SARIMA(p, d, q)(P, D, Q)}$_s$ \textbf{model} extends the ARIMA
model to handle seasonal patterns. It is defined as:
\[
\phi(B) \Phi(B^s) \nabla^d \nabla_s^D X_t = \theta(B) \Theta(B^s) Z_t
\]
or equivalently:
\[
\phi(B) \Phi(B^s) (1 - B)^d (1 - B^s)^D X_t = \theta(B) \Theta(B^s) Z_t
\]
where:
\begin{itemize}
    \item $\phi(B)$ and $\theta(B)$ are the non-seasonal AR and MA 
          polynomials.
    \item $\Phi(B^s)$ and $\Theta(B^s)$ are the seasonal AR and MA 
          polynomials.
    \item $s$ is the seasonal period (e.g., 12 for monthly data).
    \item $D$ is the order of seasonal differencing.
\end{itemize}

\subsection*{Introduction to SARIMA Processes}
The SARIMA model extends the ARIMA model to handle seasonal patterns 
in the data. It incorporates both non-seasonal and seasonal components, 
making it suitable for time series with strong seasonal trends. The 
seasonal differencing $\nabla_s^D$ removes seasonal trends, allowing 
the ARIMA model to be applied.

\subsection*{Properties}
\begin{itemize}
    \item SARIMA models are useful for time series with both trend 
          and seasonal components.
    \item The seasonal differencing $\nabla_s^D = (1 - B^s)^D$ removes 
          seasonal trends.
\end{itemize}

\section{Fractional Differencing and Long-Memory Models}
\textbf{Fractional ARIMA (ARFIMA) models} extend ARIMA models by 
allowing the differencing parameter $d$ to be a non-integer. This 
leads to \textbf{long-memory processes}, where the autocorrelation 
function decays very slowly.

\subsection*{Long-Memory Processes}
\begin{itemize}
    \item In a long-memory process, the autocorrelations are not 
          absolutely summable, meaning they decay at a hyperbolic 
          rate rather than exponentially.
    \item These processes are useful for modeling time series where 
          observations far apart in time still exhibit some correlation.
    \item The fractional differencing parameter $d$ controls the 
          degree of long memory. For $0 < d < 0.5$, the process is 
          stationary and exhibits long memory.
\end{itemize}

\section{The General Linear Process}
A \textbf{general linear process} can be expressed as an infinite-order 
MA process:
\[
X_t = \sum_{i=0}^{\infty} \psi_i Z_{t-i}
\]
where $\psi_i$ are weights and $Z_t$ is a purely random process.

\subsection*{Properties}
\begin{itemize}
    \item This representation is useful for understanding the 
          properties of stationary processes.
        \item The process is (weakly) stationary as soon as the weights are square summable,
          $\sum_i \psi_i^2 < \infty$ (then $\gamma(k) = \sigma_Z^2 \sum_i \psi_i \psi_{i+k}$). Absolute
          summability, $\sum_i |\psi_i| < \infty$, is a stronger sufficient condition that is often
          assumed because it also guarantees absolutely summable autocovariances (short memory).
\end{itemize}

\section{Continuous-Time Processes}
The chapter briefly discusses continuous-time processes, such as the 
\textbf{Ornstein-Uhlenbeck process}, which is a continuous-time 
analogue of the AR(1) process.

\subsection*{Properties}
\begin{itemize}
    \item Continuous-time processes are more complex to analyze 
          than discrete-time processes.
    \item They are often used in physics and engineering.
\end{itemize}

\section{The Wold Decomposition Theorem}
The \textbf{Wold Decomposition Theorem} states that any stationary 
process can be decomposed into a purely deterministic component 
and a purely indeterministic component, which can be represented 
as an infinite-order MA process.

\subsection*{Implications}
\begin{itemize}
    \item The theorem provides a theoretical foundation for the 
          use of MA and ARMA models in time series analysis.
    \item It highlights the importance of the purely random 
          component in modeling time series.
\end{itemize}
\end{notechapter}
